%% kapa: exam=2 round=28 year=2017 period=1 subject=appraisal_practice subject_ko=감정평가실무 seq=3 number=3 points=20
%% source: Q-Net 감정평가사 기출문제 (cst00309) artlSeq=5201712 file=raw/1008/5201712/2017년 감정평가사 2차시험 기출문제(감정평가 실무).pdf
%% conversion: kordoc 4.24.51 -> tools/gen2.py (auto); review=proofread-sonnet; flags=none
\begin{longquestion}{【문제 3】}
감정평가사 김○○은 A법원으로부터 소송 참고 용 토지 임대료 평가를 의뢰 받고, 관련 법률 등을 검토한 결과 적산법을 적용하여 관련 토지의 임대료를 평가하기로 하였다. 관련 법규 및 이론을 참작하고 제시된 자료를 활용하여 다음 물음에 답하시오. (20점) \par
\medskip\noindent 물음 1) 본건 토지의 기초가액을 구하시오. (10점) \par
\medskip\noindent 물음 2) 기대이율 결정이유를 기술하고 본건 토지의 임대료를 구하시오. (10점) \par
<자료 1> 법원 감정신청 내용 \par
1. 감정목적물 : A시 B구 C동 60-1 대 200m\textsuperscript{2} \par
2. 감정사항 : 2013.07.01부터 1년간 임대료 및 2017.07.01부터 1년간 임대료를 각각 감정평가 할 것 \par
3. 기준시점 : 2013.07.01 및 2017.07.01 \par
<자료 2> 기본적 사항 \par
1. 대상물건 \par
{\small\centering
\begin{tabular}{|c|c|c|c|c|}
\hline
\makecell{소재지} & \makecell{지목} & \makecell{면적(m\textsuperscript{2})} & \makecell{이용상황} & \makecell{용도지역} \\
\hline
\makecell{A시B구C동 60-1} & \makecell{대} & \makecell{200} & \makecell{단독주택} & \makecell{2종일주} \\
\hline
\end{tabular}\par}

2. 대상물건의 지목, 면적, 이용상황, 용도지역 내역은 2010.01.01 ～ 2017.07.01 현재까지 동일함 \par
<자료 3> 토지 개황 \par
1. 인근은 단독주택, 다세대 및 다가구 주택, 근린생활시설 등이 혼재되어 있음 \par
2. 본건은 인접 토지 대비 등고 평탄한 세장형 토지임 \par
3. 본건은 북측으로 약 6m의 포장도로에, 동측으로 약 6m의 포장도로에 접해 있음 \par
4. 토지이용계획사항 등 : 도시지역, 제2종일반주거지역 \par
5. 상기 1. ∼ 4.의 내용은 2010.01.01 ∼ 2017.07.01 현재까지 동일함 \par
<자료 4> 표준지 자료 \par
1. 인근지역내 표준지 내역(소재지 A시 B구 C동, 공시지가는 매년 1월 1일 기준) \par
{\small\centering\footnotesize\setlength{\tabcolsep}{3pt}
\begin{tabular}{|c|c|c|c|c|c|c|c|c|}
\hline
\makecell{기호} & \makecell{지번} & \makecell{지목} & \makecell{면적\newline{}(m\textsuperscript{2})} & \makecell{이용상황} & \makecell{용도지역} & \makecell{도로교통} & \makecell{형상지세} & \makecell{주위\newline{}환경} \\
\hline
\makecell{가} & \makecell{10-1} & \makecell{대} & \makecell{170} & \makecell{단독주택} & \makecell{2종일주} & \makecell{세로\newline{}(가)} & \makecell{세장형\newline{}평지} & \makecell{기존\newline{}주택지대} \\
\hline
\makecell{나} & \makecell{50-1} & \makecell{대} & \makecell{180} & \makecell{단독주택} & \makecell{2종일주} & \makecell{세로\newline{}(가)} & \makecell{세장형\newline{}평지} & \makecell{주택 및\newline{}상가지대} \\
\hline
\makecell{다} & \makecell{85-1} & \makecell{대} & \makecell{160} & \makecell{주상용} & \makecell{2종일주} & \makecell{소로\newline{}각지} & \makecell{사다리\newline{}평지} & \makecell{주택 및\newline{}상가지대} \\
\hline
\makecell{라} & \makecell{100} & \makecell{대} & \makecell{210} & \makecell{단독주택} & \makecell{2종일주} & \makecell{소로\newline{}각지} & \makecell{세장형\newline{}평지} & \makecell{주택 및\newline{}상가지대} \\
\hline
\end{tabular}\par\medskip
\begin{tabular}{|c|c|c|c|c|c|}
\hline
\multirow{2}{*}{기호} & \multicolumn{5}{c|}{공시지가(원/m\textsuperscript{2})} \\
\cline{2-6}
 & 2013년 & 2014년 & 2015년 & 2016년 & 2017년 \\
\hline
\makecell{가} & \makecell{2,830,000} & \makecell{2,950,000} & \makecell{3,150,000} & \makecell{3,270,000} & \makecell{3,430,000} \\
\hline
\makecell{나} &  &  & \makecell{2,440,000} & \makecell{2,560,000} & \makecell{2,920,000} \\
\hline
\makecell{다} & \makecell{3,000,000} & \makecell{3,230,000} & \makecell{3,350,000} & \makecell{3,540,000} & \makecell{3,700,000} \\
\hline
\makecell{라} & \makecell{2,650,000} & \makecell{2,730,000} &  &  &  \\
\hline
\end{tabular}\par}

2. 상기 1.에서 공시지가가 표시되지 않은 연도는 해당 토지가 표준지로 선정되지 않은 연도임 \par
<자료 5> 시점수정자료[지가변동률(A시 B구 주거지역, \%)] \par
{\small\centering
\begin{tabular}{|c|c|}
\hline
\makecell{구분} & \makecell{변동률} \\
\hline
\makecell{2013.01.01 ～ 2013.06.30(누계)} & \makecell{2.000} \\
\hline
\makecell{2013.07.01 ～ 2013.07.31(당월)} & \makecell{0.350} \\
\hline
\makecell{2017.01.01 ～ 2017.05.31(누계)} & \makecell{1.500} \\
\hline
\makecell{2017.05.01 ～ 2017.05.31(당월)} & \makecell{0.200} \\
\hline
\end{tabular}\par}

※ 2017년 6월 이후의 지가변동률은 현재 미고시인 상태로 직전 월인 2017년 5월 지가변동률을 연장적용 하기로 함 \par
<자료 6> 인근지역내 토지거래사례(소재지 A시 B구 C동) \par
{\small\centering\footnotesize\setlength{\tabcolsep}{3pt}
\begin{tabular}{|c|c|c|c|c|c|c|c|c|c|c|}
\hline
\makecell{기호} & \makecell{지번} & \makecell{지목} & \makecell{면적\newline{}(m\textsuperscript{2})} & \makecell{이용상황} & \makecell{용도지역} & \makecell{도로교통} & \makecell{형상\newline{}지세} & \makecell{주위환경} & \makecell{토지단가\newline{}(원/m\textsuperscript{2})} & \makecell{거래시점} \\
\hline
\makecell{1} & \makecell{20-1} & \makecell{대} & \makecell{157} & \makecell{주상용} & \makecell{2종일주} & \makecell{중로\newline{}각지} & \makecell{가장형\newline{}평지} & \makecell{주택 및\newline{}상가지대} & \makecell{4,500,000} & \makecell{2013.07.01} \\
\hline
\makecell{2} & \makecell{30-1} & \makecell{대} & \makecell{250} & \makecell{단독주택} & \makecell{2종일주} & \makecell{세로\newline{}(가)} & \makecell{세장형\newline{}평지} & \makecell{주택 및\newline{}상가지대} & \makecell{4,000,000} & \makecell{2013.07.01} \\
\hline
\makecell{3} & \makecell{70-1} & \makecell{대} & \makecell{189} & \makecell{단독주택} & \makecell{2종일주} & \makecell{소로한면} & \makecell{세장형\newline{}평지} & \makecell{기존\newline{}주택지대} & \makecell{5,100,000} & \makecell{2017.06.20} \\
\hline
\makecell{4} & \makecell{80-1} & \makecell{대} & \makecell{210} & \makecell{단독주택} & \makecell{2종일주} & \makecell{소로\newline{}각지} & \makecell{세장형\newline{}평지} & \makecell{주택 및\newline{}상가지대} & \makecell{5,400,000} & \makecell{2017.06.10} \\
\hline
\end{tabular}\par}

<자료 7> 개별요인 비교자료 \par
1. 도로접면 격차율 \par
{\small\centering
\begin{tabular}{|c|c|c|c|c|c|}
\hline
\makecell{구분} & \makecell{광대한면} & \makecell{중로한면} & \makecell{소로한면} & \makecell{세로(가)} & \makecell{세로(불)} \\
\hline
\makecell{격차율} & \makecell{1.00} & \makecell{0.90} & \makecell{0.88} & \makecell{0.85} & \makecell{0.70} \\
\hline
\end{tabular}\par}

※ 도로접면이 각지인 경우는 한면에 접하는 경우에 비해 1\% 우세함 \par
2. 형상 격차율 \par
{\small\centering
\begin{tabular}{|c|c|c|c|c|c|}
\hline
\makecell{구분} & \makecell{정방형} & \makecell{가장형} & \makecell{세장형} & \makecell{사다리} & \makecell{부정형} \\
\hline
\makecell{격차율} & \makecell{1.00} & \makecell{0.99} & \makecell{0.99} & \makecell{0.98} & \makecell{0.90} \\
\hline
\end{tabular}\par}

<자료 8> 기대이율 관련 자료 \par
1. 참고용 기대이율 적용기준율표 Ⅰ (2016년 이후 기준) \par
{\small\centering
\begin{tabular}{|c|c|c|c|c|c|}
\hline
\multicolumn{2}{|c|}{\multirow{2}{*}{\makecell{대분류}}} & \multicolumn{2}{|c|}{\multirow{2}{*}{\makecell{소분류}}} & \multicolumn{2}{|c|}{\makecell{실제이용상황}} \\
\cline{5-6}
\multicolumn{2}{|c|}{} & \multicolumn{2}{|c|}{} & \makecell{표준적이용} & \makecell{임시적이용} \\
\hline
\multirow{4}{*}{\makecell{Ⅰ}} & \multirow{4}{*}{\makecell{주거용}} & \makecell{아파트} & \multirow{4}{*}{\makecell{수도권 및\newline{}광역시}} & \makecell{1.5\%∼3.5\%} & \makecell{0.5\%∼2.5\%} \\
\cline{3-3} \cline{5-6}
 &  & \makecell{연립ㆍ다세대} &  & \makecell{1.5\%∼5.0\%} & \makecell{0.5\%∼3.0\%} \\
\cline{3-3} \cline{5-6}
 &  & \makecell{다가구} &  & \makecell{2.0\%∼6.0\%} & \makecell{1.0\%∼3.0\%} \\
\cline{3-3} \cline{5-6}
 &  & \makecell{단독주택} &  & \makecell{1.0\%∼4.0\%} & \makecell{0.5\%∼2.0\%} \\
\hline
\end{tabular}\par}

2. 참고용 기대이율 적용기준율표 Ⅱ (2016년 이전 기준) \par
{\small\centering
\begin{tabular}{|c|c|c|c|c|}
\hline
\multicolumn{2}{|c|}{\multirow{2}{*}{\makecell{토지용도}}} & \multicolumn{3}{|c|}{\makecell{실제이용상황}} \\
\cline{3-5}
\multicolumn{2}{|c|}{} & \makecell{최유효이용} & \makecell{임시적이용} & \makecell{나지} \\
\hline
\multirow{3}{*}{\makecell{주거용지}} & \makecell{아파트,연립주택,다세대주택} & \makecell{4∼7\%} & \makecell{2∼4\%} & \makecell{1∼2\%} \\
\cline{2-5}
 & \makecell{다중주택, 다가구주택} & \makecell{3∼6\%} & \makecell{2∼3\%} & \makecell{1∼2\%} \\
\cline{2-5}
 & \makecell{일반단독주택} & \makecell{3∼5\%} & \makecell{1∼3\%} & \makecell{1∼2\%} \\
\hline
\end{tabular}\par}

3. 참고용 CD금리 기준 기대이율표 (모든 년도에 적용가능) \par
{\small\centering
\begin{tabular}{|c|c|c|c|c|c|}
\hline
\multicolumn{2}{|c|}{\multirow{2}{*}{\makecell{대분류}}} & \multicolumn{2}{|c|}{\multirow{2}{*}{\makecell{소분류}}} & \multicolumn{2}{|c|}{\makecell{실제이용상황}} \\
\cline{5-6}
\multicolumn{2}{|c|}{} & \multicolumn{2}{|c|}{} & \makecell{표준적이용} & \makecell{임시적이용} \\
\hline
\multirow{4}{*}{\makecell{Ⅰ}} & \multirow{4}{*}{\makecell{주거용}} & \makecell{아파트} & \multirow{4}{*}{\makecell{수도권 및\newline{}광역시}} & \makecell{CD금리+\newline{}-1.5\%∼0.5\%} & \makecell{CD금리+\newline{}-2.5\%∼-0.5\%} \\
\cline{3-3} \cline{5-6}
 &  & \makecell{연립ㆍ다세대} &  & \makecell{CD금리+\newline{}-1.5\%∼2.0\%} & \makecell{CD금리+\newline{}-2.5\%∼ 0.0\%} \\
\cline{3-3} \cline{5-6}
 &  & \makecell{다가구} &  & \makecell{CD금리+\newline{}-1.0\%∼3.0\%} & \makecell{CD금리+\newline{}-2.0\%∼ 0.0\%} \\
\cline{3-3} \cline{5-6}
 &  & \makecell{단독주택} &  & \makecell{CD금리+\newline{}-2.0\%∼1.0\%} & \makecell{CD금리+\newline{}-2.5\%∼-1.0\%} \\
\hline
\end{tabular}\par}

4. 각종 금리 변동상황 \par
{\small\centering
\begin{tabular}{|c|c|c|}
\hline
\makecell{구분} & \makecell{국고채수익률(\%)} & \makecell{CD유통수익률(\%)} \\
\hline
\makecell{2013.07.01} & \makecell{3.10} & \makecell{3.00} \\
\hline
\makecell{2014.07.01} & \makecell{2.70} & \makecell{2.60} \\
\hline
\makecell{2015.07.01} & \makecell{2.30} & \makecell{2.20} \\
\hline
\makecell{2016.07.01} & \makecell{2.10} & \makecell{2.10} \\
\hline
\makecell{2017.07.01} & \makecell{2.00} & \makecell{2.00} \\
\hline
\end{tabular}\par}

5. 기대이율은 국고채수익률 및 CD유통수익률과 밀접한 관계가 있는 것으로 조사됨 \par
6. 상기 1. ～ 3.에서 제시하는 기대이율범위에서 각각의 중간값 중 하나를 기대이율로 선정하되, 선정된 기대이율은 2\%를 초과해야함 \par
<자료 9> 기타자료 \par
1. 토지임대료와 관련된 필요제경비는 미미하여 고려하지 않음 \par
2. 표준지 공시지가를 기준으로 토지를 평가하는 경우 그 밖의 요인 보정이 필요한 \par
것으로 판단됨 \par
3. 토지거래사례는 그 밖의 요인 보정치 산정에만 활용할 것 \par
4. 토지거래사례는 모두 정상적인 사례로 판단됨 \par
5. A시는 수도권에 위치해 있음 \par
\end{longquestion}
